% !TEX root = report.tex
% ============================================================================
% MC714 - Trabalho 2 - Report (IEEE, two columns, max. 6 pages + refs)
%
% Build:  latexmk -pdf -outdir=build report.tex
% VS Code: LaTeX Workshop compiles on save (see .vscode/settings.json)
%
% Grading rubric -> sections:
%   30% algorithm correctness + tests ..... Sec. 4, 5
%   20% model, architecture, failures ..... Sec. 2, 3, 8
%   25% experiments and reproducibility ... Sec. 6, 7
%   15% report ............................ whole document
%   10% video ............................. linked on the first page
% ============================================================================
\documentclass[conference]{IEEEtran}

\usepackage[utf8]{inputenc}
\usepackage[T1]{fontenc}
\usepackage[english]{babel}
\usepackage{microtype}
\usepackage{amsmath,amssymb}
\usepackage{graphicx}
\usepackage{booktabs}
\usepackage{url}
\usepackage{xcolor}
\usepackage[hidelinks]{hyperref}

% TODO marker, visible in the PDF; remove before submission
\newcommand{\todo}[1]{\textcolor{red}{\textbf{[TODO:} #1\textbf{]}}}

\begin{document}

\title{Raft and KV-Raft: Consensus and State Machine Replication}

\author{
\IEEEauthorblockN{Everton Romanzini Colombo}
\IEEEauthorblockA{\texttt{e257234@dac.unicamp.br}}
\and
\IEEEauthorblockN{Student 2}
\IEEEauthorblockA{\texttt{student2@unicamp.br}}
}

\maketitle

% % ---------------------------------------------------------------------------
% % Required on the first page (assignment, Section 1)
% % ---------------------------------------------------------------------------
% \noindent
% \textbf{RA:} 000000 (Student 1), 000000 (Student 2)\\
% \textbf{Repository:} \url{https://gitlab.ic.unicamp.br/.../mc714-trabalho2}\\
% \textbf{Video:} \url{https://...}

% \vspace{0.5em}

\begin{abstract}
\todo{150--200 words: problem, algorithms (leader election and
consensus/replication), containerized environment, metrics, and the main
result (linearizable under faults; latency cost under loss/delay).}
\end{abstract}

% ===========================================================================
\section{Introduction}
% ===========================================================================
\todo{Motivation: why consensus matters; why Raft (understandable) and
replicated state machines.}
\todo{Objectives: implement two complementary algorithms from the assignment
(leader election [election theme] and consensus/replication [consensus theme])
and evaluate them experimentally.}
\todo{Contributions: (i) full Raft (2A--2C) with persistence and snapshots;
(ii) linearizable KV with per-client deduplication; (iii) container packaging
with TCP RPC and an HTTP API; (iv) test suite and reproducible experiments with
linearizability checked by Porcupine.}
\todo{Paper organization (one paragraph).}

% ===========================================================================
\section{System Model}
% ===========================================================================
\todo{State and justify each item below (explicit requirement in the
assignment). Do not leave it as a bare list: justify the choices.}

\subsection{Timing model}
\todo{Partially synchronous: no known bound on delay in steady state, but
timeouts (500--700\,ms election, 150\,ms heartbeat) eventually hold, which is
what provides liveness. Justify why not purely asynchronous.}

\subsection{Channels}
\todo{Point-to-point channels over TCP (\texttt{net/rpc}), no ordering between
concurrent requests, with loss/delay injected by \texttt{tc netem};
duplication and reordering are handled by idempotence (per-client
\texttt{ClientId,Seq} dedup) and idempotent RPCs. In the test environment,
\texttt{labrpc} simulates loss, delay and partitions.}

\subsection{Failures}
\todo{Failure model: crash-stop and crash-recovery (state persisted on disk
through \texttt{Persister}); network partitions; no Byzantine faults. Justify
why persisted log plus snapshots suffice for recovery.}

\subsection{Clocks, persistence and restart}
\todo{No synchronized clocks; timeouts use the local clock only as a
heuristic. Persistence: \texttt{currentTerm}, \texttt{votedFor}, log and
\texttt{lastIncludedIndex} in \texttt{raft.state}; snapshot in \texttt{kv.snap}
(atomic write via temporary file + rename).}

\subsection{Correctness properties}
\todo{Safety: at most one leader per term; Log Matching; if an entry is applied
on one node it is applied at the same index on every node; KV linearizability.
Liveness: after the network stabilizes and with a live majority, a leader is
eventually elected and submitted commands are eventually committed.}

% ===========================================================================
\section{Architecture}
% ===========================================================================
\todo{Figure/diagram: HTTP client \(\rightarrow\) nodes (API :9000) and Raft
mesh (:8000) between containers; persistent volumes.}
\todo{Describe \texttt{compose.yaml}: 3+ nodes, bridge network, ports
9000--9002 published, 8000 internal, one volume per node.}
\todo{Communication: binary RPC (\texttt{net/rpc}/gob) between peers;
HTTP/JSON API for clients; \texttt{Clerk} on the client with retry and leader
switching. Explain why the client (not the server) finds the leader.}
\todo{Observability: \texttt{/status} endpoint (term, leader, indexes, state
and snapshot sizes) used by the experiments.}
\todo{Versions: Go, Docker/Compose, base image --- fill in
Table~\ref{tab:env}.}

% ===========================================================================
\section{Implemented Algorithms}
% ===========================================================================
\todo{Make explicit that these are the two complementary algorithms: leader
election (theme 1) and consensus/replication (theme 2).}

\subsection{Leader election and terms}
\todo{Randomized timeouts, voting (log up-to-date rule, one vote per term),
conversion to follower on a higher term; role of the observed
\texttt{leaderId}.}

\subsection{Log replication and consensus}
\todo{\texttt{AppendEntries} with \texttt{PrevLogIndex/Term}, retry with
\texttt{XTerm/XIndex/XLen}, \texttt{commitIndex} advanced only for entries from
the current term (Figure 8 of the paper), snapshots and
\texttt{InstallSnapshot}.}

\subsection{Linearizable KV state machine}
\todo{Every operation (reads included) goes through the log; response cache
keyed by \texttt{(ClientId, Seq)} for exactly-once application; replies matched
by log index; snapshot of the map plus the cache.}

\subsection{Reused code and modifications}
\todo{Mandatory citation: MIT 6.824 labs 2A--3B skeleton (\texttt{labrpc},
\texttt{labgob}, \texttt{porcupine}, tests)~\cite{mit6824}; describe what was
changed: the Raft/KV implementation, the \texttt{RPCClient} interface for real
TCP, the on-disk \texttt{Persister}, the \texttt{/status} endpoint, the
benchmark client and the experiment scripts.}

% ===========================================================================
\section{Correctness Properties and Tests}
% ===========================================================================
\subsection{Safety arguments}
\todo{Sketch the proof: leader uniqueness per term; majority intersection;
the election restriction guarantees the elected leader holds every committed
entry; in-order application by index.}

\subsection{Automated tests}
\todo{Describe the inherited suite (2A/2B/2C, 3A/3B) and our own tests: real
TCP cluster, follower rejecting writes (HTTP 421), \texttt{/status}
consistency. Report that everything passes (total time).}
\todo{Linearizability: Porcupine reconstructs a valid order for the history;
46/46 runs \texttt{linearizable=ok} in the experimental scenarios.}

% ===========================================================================
\section{Experimental Methodology}
% ===========================================================================
\begin{table}[t]
\centering
\small
\caption{Execution environment}
\label{tab:env}
\begin{tabular}{ll}
\toprule
Item & Value \\
\midrule
Hardware & \todo{CPU, cores, RAM} \\
Operating system & \todo{distribution, kernel} \\
Docker / Compose & \todo{versions} \\
Go & \todo{version} \\
Nodes & 3, 5, 7 containers (same host) \\
Load & 8 clients, 50 ops/client \\
Seeds & 5--10 per configuration \\
Warm-up & 1 discarded run (sweep) \\
\bottomrule
\end{tabular}
\end{table}

\todo{Describe the workload (50/50 read/write mix, 5 keys, \texttt{kvbench}),
the per-run key namespace, and the fact that throughput is measured before the
linearizability check.}
\todo{Metrics: throughput (ops/s), latency p50/p95/p99, time-to-elect,
service recovery time (TTR-service), replica recovery time (TTR-replica),
attempts per operation and \texttt{ErrWrongLeader} rate. Justify each one.}
\todo{Parameters varied at \(\geq 3\) levels: \(N \in \{3,5,7\}\); delay/loss
via \texttt{tc netem} (0, 50\,ms, 100\,ms; 0\%, 2\%, 5\%); number of clients.}
\todo{Fault scenarios: leader crash, leader partition, full cluster restart.}
\todo{Reproducibility: \texttt{scripts/experiment.py} with recorded seeds,
raw CSVs in \texttt{scripts/experiments/}, analysis in
\texttt{scripts/analysis.ipynb}. Exact commands.}

% ===========================================================================
\section{Results}
% ===========================================================================
\todo{Replace the values below with the final version and generate the figures
with \texttt{scripts/analysis.ipynb}.}

\begin{table}[t]
\centering
\small
\caption{Throughput and latency (median over seeds)}
\label{tab:res}
\setlength{\tabcolsep}{4pt}
\begin{tabular}{lrrrr}
\toprule
Scenario & ops/s & p50 & p95 & p99 \\
\midrule
Normal            & 1633 & 4.5  & 7.4  & 8.8 \\
\texttt{netem} 50\,ms/2\% & 41 & 155 & 307 & 418 \\
Leader crash (2 nodes) & 2038 & 2.7 & 8.0 & 21.9 \\
Partition (after heal) & 1227 & 2.6 & 6.4 & 184 \\
\bottomrule
\end{tabular}
\end{table}

\begin{table}[t]
\centering
\small
\caption{Recovery times (mean \(\pm\) std)}
\label{tab:ttr}
\setlength{\tabcolsep}{4pt}
\begin{tabular}{lrrr}
\toprule
Scenario & election & TTR-service & TTR-replica \\
\midrule
Leader crash & \(507 \pm 114\)\,ms & \(1868 \pm 32\)\,ms & \(468 \pm 124\)\,ms \\
Partition    & \(349 \pm 185\)\,ms & \(1632 \pm 404\)\,ms & not applicable \\
\bottomrule
\end{tabular}
\end{table}

% Suggested figures (export from scripts/analysis.ipynb to figs/):
% \begin{figure}[t]
% \centering
% \includegraphics[width=\columnwidth]{figs/throughput-vs-n.png}
% \caption{Mean throughput with \(N \in \{3,5,7\}\) (bars: standard deviation).}
% \label{fig:n}
% \end{figure}
%
% \begin{figure}[t]
% \centering
% \includegraphics[width=\columnwidth]{figs/probe-latency.png}
% \caption{Write latency during the leader crash; dashed line: fault instant.}
% \label{fig:probe}
% \end{figure}

\subsection{Normal operation}
\todo{Describe throughput and percentiles; explain that the 2\% retry in the
baseline is one extra attempt per client to discover the leader (8/400), not
steady-state retrying.}

\subsection{Delay and loss}
\todo{Discussion: throughput drop and p99 increase; why the ceiling is
dominated by the election timeout and retransmissions; linearizability is
preserved.}

\subsection{Leader crash}
\todo{Time-to-elect \(\approx\) election timeout; TTR-service dominated by the
client HTTP timeout (\(\sim\)2\,s); TTR-replica measures recovery via
log/snapshot.}

\subsection{Partition and the disruptive-server problem}
\todo{Inflated terms, churn, attempts per operation; explain that without
PreVote the isolated node exiles itself; treat it as a limitation and future
work.}

\subsection{Scalability}
\todo{Results with \(N=3,5,7\): throughput and latency; discuss why the leader
serializes and where replication cost would show up (message counts).}

% ===========================================================================
\section{Limitations and Threats to Validity}
% ===========================================================================
\todo{Containers on one host are not independent failure domains; variability
from short runs; no batching; no PreVote; no message/byte counters; Docker
Desktop environment.}

% ===========================================================================
\section{Conclusion and Future Work}
% ===========================================================================
\todo{Recap: two algorithms, correctness under faults, performance cost.
Future work: PreVote/check-quorum, batching, lease-based reads, incremental
snapshots, message accounting, execution across multiple machines.}

% ===========================================================================
\bibliographystyle{IEEEtran}
\bibliography{refs}

\end{document}
